% ============================================================
% DAFTAR PUSTAKA — SCOPE: CAN (FRONTEND)
% Farchan Deano Muhammad — D4 TRM PENS
% Proyek: Sketchbook Universe — Interactive HITL AI Literacy Simulation
%
% Total: 32 referensi (20 General/Shared + 12 Can-specific)
% Format: IEEE numbered, urut by appearance (Bab 1 -> Bab 3)
%
% Source: Sitasi_Final_v1.md (16 Juni 2026)
% Catatan:
%   - ref10 (B02 Meng et al. MediaPipe) tetap dimasukkan dengan catatan
%     DEPRECATED post-pivot 16/6/26 (login gesture diganti nomor absen).
%   - ref10 dipertahankan untuk dokumentasi historis canvas drawing input.
%
% Cara pakai:
%   % ============================================================
% DAFTAR PUSTAKA — SCOPE: CAN (FRONTEND)
% Farchan Deano Muhammad — D4 TRM PENS
% Proyek: Sketchbook Universe — Interactive HITL AI Literacy Simulation
%
% Total: 32 referensi (20 General/Shared + 12 Can-specific)
% Format: IEEE numbered, urut by appearance (Bab 1 -> Bab 3)
%
% Source: Sitasi_Final_v1.md (16 Juni 2026)
% Catatan:
%   - ref10 (B02 Meng et al. MediaPipe) tetap dimasukkan dengan catatan
%     DEPRECATED post-pivot 16/6/26 (login gesture diganti nomor absen).
%   - ref10 dipertahankan untuk dokumentasi historis canvas drawing input.
%
% Cara pakai:
%   % ============================================================
% DAFTAR PUSTAKA — SCOPE: CAN (FRONTEND)
% Farchan Deano Muhammad — D4 TRM PENS
% Proyek: Sketchbook Universe — Interactive HITL AI Literacy Simulation
%
% Total: 32 referensi (20 General/Shared + 12 Can-specific)
% Format: IEEE numbered, urut by appearance (Bab 1 -> Bab 3)
%
% Source: Sitasi_Final_v1.md (16 Juni 2026)
% Catatan:
%   - ref10 (B02 Meng et al. MediaPipe) tetap dimasukkan dengan catatan
%     DEPRECATED post-pivot 16/6/26 (login gesture diganti nomor absen).
%   - ref10 dipertahankan untuk dokumentasi historis canvas drawing input.
%
% Cara pakai:
%   % ============================================================
% DAFTAR PUSTAKA — SCOPE: CAN (FRONTEND)
% Farchan Deano Muhammad — D4 TRM PENS
% Proyek: Sketchbook Universe — Interactive HITL AI Literacy Simulation
%
% Total: 32 referensi (20 General/Shared + 12 Can-specific)
% Format: IEEE numbered, urut by appearance (Bab 1 -> Bab 3)
%
% Source: Sitasi_Final_v1.md (16 Juni 2026)
% Catatan:
%   - ref10 (B02 Meng et al. MediaPipe) tetap dimasukkan dengan catatan
%     DEPRECATED post-pivot 16/6/26 (login gesture diganti nomor absen).
%   - ref10 dipertahankan untuk dokumentasi historis canvas drawing input.
%
% Cara pakai:
%   \input{daftar_pustaka_can_v1.tex}
%   atau copy bagian thebibliography ke proposal master.
% ============================================================

\chapter*{DAFTAR PUSTAKA}
\addcontentsline{toc}{chapter}{DAFTAR PUSTAKA}

\begin{thebibliography}{99}

\bibitem{ref1}
Ng, D.T.K., Leung, J.K.L., Chu, S.K.W., \& Qiao, M.S. (2021). Conceptualizing AI literacy: An exploratory review. \textit{Computers and Education: Artificial Intelligence}, 2, 100041.

\bibitem{ref2}
Casal-Otero, V., Catala, A., Fern\'andez-Morante, C., Taboada, M., Caeiro-Rodr\'iguez, M., \& Barro, S. (2023). AI literacy in K-12: A systematic literature review. \textit{International Journal of STEM Education}, 10, 29.

\bibitem{ref3}
Khosravi, H., Shum, S.B., Chen, G., Conati, C., Tsai, Y.-S., Kay, J., Knight, S., Martinez-Maldonado, R., Sadat-Milani, L., \& Ga\v{s}evi\'c, D. (2022). Explainable Artificial Intelligence in education. \textit{Computers and Education: Artificial Intelligence}, 3, 100074.

\bibitem{ref4}
Mosqueira-Rey, E., Hern\'andez-Pereira, E., Alonso-R\'ios, D., Bobillo-Ares, B., \& Moret-Bonillo, V. (2023). Human-in-the-loop machine learning: A state of the art. \textit{Artificial Intelligence Review}, 56(4), 3005--3054.

\bibitem{ref5}
Memarian, B., \& Doleck, T. (2024). Human-in-the-loop in artificial intelligence in education: A review and entity-relationship (ER) analysis. \textit{Computers in Human Behavior: Artificial Humans}, 2, 100053.

\bibitem{ref6}
Videnovik, M., Vlahovi\v{c}-Steti\'c, V., Kišiček, S., \& Vošner, H.B. (2023). Game-based learning in computer science education: A scoping literature review. \textit{International Journal of STEM Education}, 10, 54.

\bibitem{ref7}
Chan, K.K., Wan, K.M., \& King, R.B. (2021). Performance over enjoyment? Effect of game-based learning on learning outcome and flow experience. \textit{Frontiers in Education}, 6, 660376.

\bibitem{ref8}
Karran, A.J., Demazure, T., \& Hudelot, C. (2022). Designing for confidence: The impact of visualizing artificial intelligence decisions. \textit{Frontiers in Neuroscience}, 16.

% --- Can-specific (Kategori B) ---

\bibitem{ref9}
Schroeder, N.L., Davis, C.S., \& Yang, Y. (2025). Designing and learning with pedagogical agents: An umbrella review. \textit{Journal of Educational Computing Research}, 62(8).

\bibitem{ref10}
Meng, Z., Fu, Z., Liu, J., Wang, H., \& Pan, Y. (2024). Real-Time Hand Gesture Monitoring Model Based on MediaPipe's Registerable System. \textit{Sensors}, 24(19), 6262. \textit{[Catatan: DEPRECATED post-pivot 16/6/26 --- login gesture diganti nomor absen + superadmin. Dipertahankan untuk dokumentasi historis canvas drawing input.]}

\bibitem{ref11}
Wang, N., \& Lester, J. (2023). K-12 education in the age of AI: A call to action for K-12 AI literacy. \textit{International Journal of Artificial Intelligence in Education}, 33(3), 631--640.

\bibitem{ref12}
Yim, I. H. Y., \& Su, J. (2025). Artificial intelligence literacy education in primary schools: A review. \textit{International Journal of Technology and Design Education}, 35, Article 57.

\bibitem{ref13}
Lee, S., Mott, B., Ottenbreit-Leftwich, A., Scribner, A., Taylor, S., Glazewski, K., \& Lester, J. (2021). AI-infused collaborative inquiry in upper elementary school: A game-based learning approach. \textit{Proceedings of the AAAI Conference on Artificial Intelligence}, 35(17), 15653--15661.

\bibitem{ref14}
Gupta, A., Lee, S., Mott, B., Chakraburty, S., Scribner, A., Glazewski, K., Ottenbreit-Leftwich, A., \& Lester, J. (2024). Supporting upper elementary students in learning AI concepts with story-driven game-based learning. \textit{Proceedings of the AAAI Conference on Artificial Intelligence}, 38(21), 23321--23329.

\bibitem{ref15}
Hsu, T. C., Abelson, H., Lao, N., Tseng, Y. H., \& Lin, Y. T. (2021). Behavioral-pattern exploration and development of an instructional tool for young children to learn AI. \textit{Computers and Education: Artificial Intelligence}, 2, Article 100060.

\bibitem{ref16}
Ocak, C., Kopcha, T. J., \& Dey, R. (2023). An AI-enhanced pattern recognition approach to temporal and spatial analysis of children's embodied interactions. \textit{Computers and Education: Artificial Intelligence}, 4, Article 100155.

\bibitem{ref17}
Shokeen, E., Katirci, N., Williams-Pierce, C., \& Fan, C. (2022). Children learning to sketch: Sketching to learn. \textit{Information and Learning Sciences}, 123(7/8), 482--497.

\bibitem{ref18}
Chen, Z., Duan, Z., Qi, Y., \& Wang, W. (2025). From line sketches to contextualized scenes: An AR scaffold for narrative drawing based on sketch recognition and its motivational effects. \textit{Proceedings of the 2025 ACM/IEEE International Conference on Human-Robot Interaction (HRI)}, 1--10.

\bibitem{ref19}
Wang, C. J., Zhong, H. X., Chiu, P. S., Chang, J. H., \& Chen, M. Y. (2022). Research on the impacts of cognitive style and computational thinking on college students in a visual artificial intelligence course. \textit{Frontiers in Psychology}, 13, Article 864416.

\bibitem{ref20}
Wu, S. Y., \& Su, Y. S. (2021). Visual programming environments and computational thinking performance of fifth-and sixth-grade students. \textit{Journal of Educational Computing Research}, 59(7), 1413--1441.

% --- General/Shared additions (Kategori A baru) ---

\bibitem{ref21}
Yin, M., Wortman Vaughan, J., \& Wallach, H. (2019). Understanding the effect of accuracy on trust in machine learning models. \textit{Proceedings of the 2019 CHI Conference on Human Factors in Computing Systems}, 1--12.

\bibitem{ref22}
Batanero, C., \& \'Alvarez-Arroyo, R. (2023). Teaching and learning of probability. \textit{ZDM -- Mathematics Education}, 56, 5--17.

\bibitem{ref23}
Zhan, Z., Tong, Y., Lan, X., \& Zhong, B. (2024). A systematic literature review of game-based learning in artificial intelligence education. \textit{Interactive Learning Environments}, 32(4), 1189--1212.

\bibitem{ref24}
Liao, Q. V., Gruen, D., \& Miller, S. (2020). Questioning the AI: Informing design practices for explainable AI user experiences. \textit{Proceedings of the 2020 CHI Conference on Human Factors in Computing Systems}, 1--15.

\bibitem{ref25}
Feldman-Maggor, Y., Cukurova, M., Kent, C., Luckin, R., \& Baur-Wolfson, M. (2025). The impact of explainable AI on teachers' trust and acceptance of AI EdTech recommendations. \textit{International Journal of Artificial Intelligence in Education}, 35(1), 1--25.

% --- Methodology (Kategori A existing, methodology) ---

\bibitem{ref26}
Pressman, R.S. (2014). \textit{Software Engineering: A Practitioner's Approach} (8th ed.). New York: McGraw-Hill.

\bibitem{ref27}
Sommerville, I. (2015). \textit{Software Engineering} (10th ed.). Pearson.

\bibitem{ref28}
Sugiyono. (2020). \textit{Metode Penelitian Kuantitatif, Kualitatif, dan R\&D}. Bandung: Alfabeta.

\bibitem{ref29}
Creswell, J.W. (2014). \textit{Research Design: Qualitative, Quantitative, and Mixed Methods Approaches} (4th ed.). Sage Publications.

\bibitem{ref30}
Ishikawa, K. (1986). \textit{Guide to Quality Control} (2nd ed.). Tokyo: Asian Productivity Organization.

\bibitem{ref31}
Brooke, J. (1996). SUS: A `Quick and Dirty' Usability Scale. In \textit{Usability Evaluation in Industry} (pp. 189--194). Taylor \& Francis.

\bibitem{ref32}
Bangor, A., Kortum, P.T., \& Miller, J.T. (2008). An empirical evaluation of the System Usability Scale. \textit{International Journal of Human-Computer Interaction}, 24(6), 574--594.

\end{thebibliography}

%   atau copy bagian thebibliography ke proposal master.
% ============================================================

\chapter*{DAFTAR PUSTAKA}
\addcontentsline{toc}{chapter}{DAFTAR PUSTAKA}

\begin{thebibliography}{99}

\bibitem{ref1}
Ng, D.T.K., Leung, J.K.L., Chu, S.K.W., \& Qiao, M.S. (2021). Conceptualizing AI literacy: An exploratory review. \textit{Computers and Education: Artificial Intelligence}, 2, 100041.

\bibitem{ref2}
Casal-Otero, V., Catala, A., Fern\'andez-Morante, C., Taboada, M., Caeiro-Rodr\'iguez, M., \& Barro, S. (2023). AI literacy in K-12: A systematic literature review. \textit{International Journal of STEM Education}, 10, 29.

\bibitem{ref3}
Khosravi, H., Shum, S.B., Chen, G., Conati, C., Tsai, Y.-S., Kay, J., Knight, S., Martinez-Maldonado, R., Sadat-Milani, L., \& Ga\v{s}evi\'c, D. (2022). Explainable Artificial Intelligence in education. \textit{Computers and Education: Artificial Intelligence}, 3, 100074.

\bibitem{ref4}
Mosqueira-Rey, E., Hern\'andez-Pereira, E., Alonso-R\'ios, D., Bobillo-Ares, B., \& Moret-Bonillo, V. (2023). Human-in-the-loop machine learning: A state of the art. \textit{Artificial Intelligence Review}, 56(4), 3005--3054.

\bibitem{ref5}
Memarian, B., \& Doleck, T. (2024). Human-in-the-loop in artificial intelligence in education: A review and entity-relationship (ER) analysis. \textit{Computers in Human Behavior: Artificial Humans}, 2, 100053.

\bibitem{ref6}
Videnovik, M., Vlahovi\v{c}-Steti\'c, V., Kišiček, S., \& Vošner, H.B. (2023). Game-based learning in computer science education: A scoping literature review. \textit{International Journal of STEM Education}, 10, 54.

\bibitem{ref7}
Chan, K.K., Wan, K.M., \& King, R.B. (2021). Performance over enjoyment? Effect of game-based learning on learning outcome and flow experience. \textit{Frontiers in Education}, 6, 660376.

\bibitem{ref8}
Karran, A.J., Demazure, T., \& Hudelot, C. (2022). Designing for confidence: The impact of visualizing artificial intelligence decisions. \textit{Frontiers in Neuroscience}, 16.

% --- Can-specific (Kategori B) ---

\bibitem{ref9}
Schroeder, N.L., Davis, C.S., \& Yang, Y. (2025). Designing and learning with pedagogical agents: An umbrella review. \textit{Journal of Educational Computing Research}, 62(8).

\bibitem{ref10}
Meng, Z., Fu, Z., Liu, J., Wang, H., \& Pan, Y. (2024). Real-Time Hand Gesture Monitoring Model Based on MediaPipe's Registerable System. \textit{Sensors}, 24(19), 6262. \textit{[Catatan: DEPRECATED post-pivot 16/6/26 --- login gesture diganti nomor absen + superadmin. Dipertahankan untuk dokumentasi historis canvas drawing input.]}

\bibitem{ref11}
Wang, N., \& Lester, J. (2023). K-12 education in the age of AI: A call to action for K-12 AI literacy. \textit{International Journal of Artificial Intelligence in Education}, 33(3), 631--640.

\bibitem{ref12}
Yim, I. H. Y., \& Su, J. (2025). Artificial intelligence literacy education in primary schools: A review. \textit{International Journal of Technology and Design Education}, 35, Article 57.

\bibitem{ref13}
Lee, S., Mott, B., Ottenbreit-Leftwich, A., Scribner, A., Taylor, S., Glazewski, K., \& Lester, J. (2021). AI-infused collaborative inquiry in upper elementary school: A game-based learning approach. \textit{Proceedings of the AAAI Conference on Artificial Intelligence}, 35(17), 15653--15661.

\bibitem{ref14}
Gupta, A., Lee, S., Mott, B., Chakraburty, S., Scribner, A., Glazewski, K., Ottenbreit-Leftwich, A., \& Lester, J. (2024). Supporting upper elementary students in learning AI concepts with story-driven game-based learning. \textit{Proceedings of the AAAI Conference on Artificial Intelligence}, 38(21), 23321--23329.

\bibitem{ref15}
Hsu, T. C., Abelson, H., Lao, N., Tseng, Y. H., \& Lin, Y. T. (2021). Behavioral-pattern exploration and development of an instructional tool for young children to learn AI. \textit{Computers and Education: Artificial Intelligence}, 2, Article 100060.

\bibitem{ref16}
Ocak, C., Kopcha, T. J., \& Dey, R. (2023). An AI-enhanced pattern recognition approach to temporal and spatial analysis of children's embodied interactions. \textit{Computers and Education: Artificial Intelligence}, 4, Article 100155.

\bibitem{ref17}
Shokeen, E., Katirci, N., Williams-Pierce, C., \& Fan, C. (2022). Children learning to sketch: Sketching to learn. \textit{Information and Learning Sciences}, 123(7/8), 482--497.

\bibitem{ref18}
Chen, Z., Duan, Z., Qi, Y., \& Wang, W. (2025). From line sketches to contextualized scenes: An AR scaffold for narrative drawing based on sketch recognition and its motivational effects. \textit{Proceedings of the 2025 ACM/IEEE International Conference on Human-Robot Interaction (HRI)}, 1--10.

\bibitem{ref19}
Wang, C. J., Zhong, H. X., Chiu, P. S., Chang, J. H., \& Chen, M. Y. (2022). Research on the impacts of cognitive style and computational thinking on college students in a visual artificial intelligence course. \textit{Frontiers in Psychology}, 13, Article 864416.

\bibitem{ref20}
Wu, S. Y., \& Su, Y. S. (2021). Visual programming environments and computational thinking performance of fifth-and sixth-grade students. \textit{Journal of Educational Computing Research}, 59(7), 1413--1441.

% --- General/Shared additions (Kategori A baru) ---

\bibitem{ref21}
Yin, M., Wortman Vaughan, J., \& Wallach, H. (2019). Understanding the effect of accuracy on trust in machine learning models. \textit{Proceedings of the 2019 CHI Conference on Human Factors in Computing Systems}, 1--12.

\bibitem{ref22}
Batanero, C., \& \'Alvarez-Arroyo, R. (2023). Teaching and learning of probability. \textit{ZDM -- Mathematics Education}, 56, 5--17.

\bibitem{ref23}
Zhan, Z., Tong, Y., Lan, X., \& Zhong, B. (2024). A systematic literature review of game-based learning in artificial intelligence education. \textit{Interactive Learning Environments}, 32(4), 1189--1212.

\bibitem{ref24}
Liao, Q. V., Gruen, D., \& Miller, S. (2020). Questioning the AI: Informing design practices for explainable AI user experiences. \textit{Proceedings of the 2020 CHI Conference on Human Factors in Computing Systems}, 1--15.

\bibitem{ref25}
Feldman-Maggor, Y., Cukurova, M., Kent, C., Luckin, R., \& Baur-Wolfson, M. (2025). The impact of explainable AI on teachers' trust and acceptance of AI EdTech recommendations. \textit{International Journal of Artificial Intelligence in Education}, 35(1), 1--25.

% --- Methodology (Kategori A existing, methodology) ---

\bibitem{ref26}
Pressman, R.S. (2014). \textit{Software Engineering: A Practitioner's Approach} (8th ed.). New York: McGraw-Hill.

\bibitem{ref27}
Sommerville, I. (2015). \textit{Software Engineering} (10th ed.). Pearson.

\bibitem{ref28}
Sugiyono. (2020). \textit{Metode Penelitian Kuantitatif, Kualitatif, dan R\&D}. Bandung: Alfabeta.

\bibitem{ref29}
Creswell, J.W. (2014). \textit{Research Design: Qualitative, Quantitative, and Mixed Methods Approaches} (4th ed.). Sage Publications.

\bibitem{ref30}
Ishikawa, K. (1986). \textit{Guide to Quality Control} (2nd ed.). Tokyo: Asian Productivity Organization.

\bibitem{ref31}
Brooke, J. (1996). SUS: A `Quick and Dirty' Usability Scale. In \textit{Usability Evaluation in Industry} (pp. 189--194). Taylor \& Francis.

\bibitem{ref32}
Bangor, A., Kortum, P.T., \& Miller, J.T. (2008). An empirical evaluation of the System Usability Scale. \textit{International Journal of Human-Computer Interaction}, 24(6), 574--594.

\end{thebibliography}

%   atau copy bagian thebibliography ke proposal master.
% ============================================================

\chapter*{DAFTAR PUSTAKA}
\addcontentsline{toc}{chapter}{DAFTAR PUSTAKA}

\begin{thebibliography}{99}

\bibitem{ref1}
Ng, D.T.K., Leung, J.K.L., Chu, S.K.W., \& Qiao, M.S. (2021). Conceptualizing AI literacy: An exploratory review. \textit{Computers and Education: Artificial Intelligence}, 2, 100041.

\bibitem{ref2}
Casal-Otero, V., Catala, A., Fern\'andez-Morante, C., Taboada, M., Caeiro-Rodr\'iguez, M., \& Barro, S. (2023). AI literacy in K-12: A systematic literature review. \textit{International Journal of STEM Education}, 10, 29.

\bibitem{ref3}
Khosravi, H., Shum, S.B., Chen, G., Conati, C., Tsai, Y.-S., Kay, J., Knight, S., Martinez-Maldonado, R., Sadat-Milani, L., \& Ga\v{s}evi\'c, D. (2022). Explainable Artificial Intelligence in education. \textit{Computers and Education: Artificial Intelligence}, 3, 100074.

\bibitem{ref4}
Mosqueira-Rey, E., Hern\'andez-Pereira, E., Alonso-R\'ios, D., Bobillo-Ares, B., \& Moret-Bonillo, V. (2023). Human-in-the-loop machine learning: A state of the art. \textit{Artificial Intelligence Review}, 56(4), 3005--3054.

\bibitem{ref5}
Memarian, B., \& Doleck, T. (2024). Human-in-the-loop in artificial intelligence in education: A review and entity-relationship (ER) analysis. \textit{Computers in Human Behavior: Artificial Humans}, 2, 100053.

\bibitem{ref6}
Videnovik, M., Vlahovi\v{c}-Steti\'c, V., Kišiček, S., \& Vošner, H.B. (2023). Game-based learning in computer science education: A scoping literature review. \textit{International Journal of STEM Education}, 10, 54.

\bibitem{ref7}
Chan, K.K., Wan, K.M., \& King, R.B. (2021). Performance over enjoyment? Effect of game-based learning on learning outcome and flow experience. \textit{Frontiers in Education}, 6, 660376.

\bibitem{ref8}
Karran, A.J., Demazure, T., \& Hudelot, C. (2022). Designing for confidence: The impact of visualizing artificial intelligence decisions. \textit{Frontiers in Neuroscience}, 16.

% --- Can-specific (Kategori B) ---

\bibitem{ref9}
Schroeder, N.L., Davis, C.S., \& Yang, Y. (2025). Designing and learning with pedagogical agents: An umbrella review. \textit{Journal of Educational Computing Research}, 62(8).

\bibitem{ref10}
Meng, Z., Fu, Z., Liu, J., Wang, H., \& Pan, Y. (2024). Real-Time Hand Gesture Monitoring Model Based on MediaPipe's Registerable System. \textit{Sensors}, 24(19), 6262. \textit{[Catatan: DEPRECATED post-pivot 16/6/26 --- login gesture diganti nomor absen + superadmin. Dipertahankan untuk dokumentasi historis canvas drawing input.]}

\bibitem{ref11}
Wang, N., \& Lester, J. (2023). K-12 education in the age of AI: A call to action for K-12 AI literacy. \textit{International Journal of Artificial Intelligence in Education}, 33(3), 631--640.

\bibitem{ref12}
Yim, I. H. Y., \& Su, J. (2025). Artificial intelligence literacy education in primary schools: A review. \textit{International Journal of Technology and Design Education}, 35, Article 57.

\bibitem{ref13}
Lee, S., Mott, B., Ottenbreit-Leftwich, A., Scribner, A., Taylor, S., Glazewski, K., \& Lester, J. (2021). AI-infused collaborative inquiry in upper elementary school: A game-based learning approach. \textit{Proceedings of the AAAI Conference on Artificial Intelligence}, 35(17), 15653--15661.

\bibitem{ref14}
Gupta, A., Lee, S., Mott, B., Chakraburty, S., Scribner, A., Glazewski, K., Ottenbreit-Leftwich, A., \& Lester, J. (2024). Supporting upper elementary students in learning AI concepts with story-driven game-based learning. \textit{Proceedings of the AAAI Conference on Artificial Intelligence}, 38(21), 23321--23329.

\bibitem{ref15}
Hsu, T. C., Abelson, H., Lao, N., Tseng, Y. H., \& Lin, Y. T. (2021). Behavioral-pattern exploration and development of an instructional tool for young children to learn AI. \textit{Computers and Education: Artificial Intelligence}, 2, Article 100060.

\bibitem{ref16}
Ocak, C., Kopcha, T. J., \& Dey, R. (2023). An AI-enhanced pattern recognition approach to temporal and spatial analysis of children's embodied interactions. \textit{Computers and Education: Artificial Intelligence}, 4, Article 100155.

\bibitem{ref17}
Shokeen, E., Katirci, N., Williams-Pierce, C., \& Fan, C. (2022). Children learning to sketch: Sketching to learn. \textit{Information and Learning Sciences}, 123(7/8), 482--497.

\bibitem{ref18}
Chen, Z., Duan, Z., Qi, Y., \& Wang, W. (2025). From line sketches to contextualized scenes: An AR scaffold for narrative drawing based on sketch recognition and its motivational effects. \textit{Proceedings of the 2025 ACM/IEEE International Conference on Human-Robot Interaction (HRI)}, 1--10.

\bibitem{ref19}
Wang, C. J., Zhong, H. X., Chiu, P. S., Chang, J. H., \& Chen, M. Y. (2022). Research on the impacts of cognitive style and computational thinking on college students in a visual artificial intelligence course. \textit{Frontiers in Psychology}, 13, Article 864416.

\bibitem{ref20}
Wu, S. Y., \& Su, Y. S. (2021). Visual programming environments and computational thinking performance of fifth-and sixth-grade students. \textit{Journal of Educational Computing Research}, 59(7), 1413--1441.

% --- General/Shared additions (Kategori A baru) ---

\bibitem{ref21}
Yin, M., Wortman Vaughan, J., \& Wallach, H. (2019). Understanding the effect of accuracy on trust in machine learning models. \textit{Proceedings of the 2019 CHI Conference on Human Factors in Computing Systems}, 1--12.

\bibitem{ref22}
Batanero, C., \& \'Alvarez-Arroyo, R. (2023). Teaching and learning of probability. \textit{ZDM -- Mathematics Education}, 56, 5--17.

\bibitem{ref23}
Zhan, Z., Tong, Y., Lan, X., \& Zhong, B. (2024). A systematic literature review of game-based learning in artificial intelligence education. \textit{Interactive Learning Environments}, 32(4), 1189--1212.

\bibitem{ref24}
Liao, Q. V., Gruen, D., \& Miller, S. (2020). Questioning the AI: Informing design practices for explainable AI user experiences. \textit{Proceedings of the 2020 CHI Conference on Human Factors in Computing Systems}, 1--15.

\bibitem{ref25}
Feldman-Maggor, Y., Cukurova, M., Kent, C., Luckin, R., \& Baur-Wolfson, M. (2025). The impact of explainable AI on teachers' trust and acceptance of AI EdTech recommendations. \textit{International Journal of Artificial Intelligence in Education}, 35(1), 1--25.

% --- Methodology (Kategori A existing, methodology) ---

\bibitem{ref26}
Pressman, R.S. (2014). \textit{Software Engineering: A Practitioner's Approach} (8th ed.). New York: McGraw-Hill.

\bibitem{ref27}
Sommerville, I. (2015). \textit{Software Engineering} (10th ed.). Pearson.

\bibitem{ref28}
Sugiyono. (2020). \textit{Metode Penelitian Kuantitatif, Kualitatif, dan R\&D}. Bandung: Alfabeta.

\bibitem{ref29}
Creswell, J.W. (2014). \textit{Research Design: Qualitative, Quantitative, and Mixed Methods Approaches} (4th ed.). Sage Publications.

\bibitem{ref30}
Ishikawa, K. (1986). \textit{Guide to Quality Control} (2nd ed.). Tokyo: Asian Productivity Organization.

\bibitem{ref31}
Brooke, J. (1996). SUS: A `Quick and Dirty' Usability Scale. In \textit{Usability Evaluation in Industry} (pp. 189--194). Taylor \& Francis.

\bibitem{ref32}
Bangor, A., Kortum, P.T., \& Miller, J.T. (2008). An empirical evaluation of the System Usability Scale. \textit{International Journal of Human-Computer Interaction}, 24(6), 574--594.

\end{thebibliography}

%   atau copy bagian thebibliography ke proposal master.
% ============================================================

\chapter*{DAFTAR PUSTAKA}
\addcontentsline{toc}{chapter}{DAFTAR PUSTAKA}

\begin{thebibliography}{99}

\bibitem{ref1}
Ng, D.T.K., Leung, J.K.L., Chu, S.K.W., \& Qiao, M.S. (2021). Conceptualizing AI literacy: An exploratory review. \textit{Computers and Education: Artificial Intelligence}, 2, 100041.

\bibitem{ref2}
Casal-Otero, V., Catala, A., Fern\'andez-Morante, C., Taboada, M., Caeiro-Rodr\'iguez, M., \& Barro, S. (2023). AI literacy in K-12: A systematic literature review. \textit{International Journal of STEM Education}, 10, 29.

\bibitem{ref3}
Khosravi, H., Shum, S.B., Chen, G., Conati, C., Tsai, Y.-S., Kay, J., Knight, S., Martinez-Maldonado, R., Sadat-Milani, L., \& Ga\v{s}evi\'c, D. (2022). Explainable Artificial Intelligence in education. \textit{Computers and Education: Artificial Intelligence}, 3, 100074.

\bibitem{ref4}
Mosqueira-Rey, E., Hern\'andez-Pereira, E., Alonso-R\'ios, D., Bobillo-Ares, B., \& Moret-Bonillo, V. (2023). Human-in-the-loop machine learning: A state of the art. \textit{Artificial Intelligence Review}, 56(4), 3005--3054.

\bibitem{ref5}
Memarian, B., \& Doleck, T. (2024). Human-in-the-loop in artificial intelligence in education: A review and entity-relationship (ER) analysis. \textit{Computers in Human Behavior: Artificial Humans}, 2, 100053.

\bibitem{ref6}
Videnovik, M., Vlahovi\v{c}-Steti\'c, V., Kišiček, S., \& Vošner, H.B. (2023). Game-based learning in computer science education: A scoping literature review. \textit{International Journal of STEM Education}, 10, 54.

\bibitem{ref7}
Chan, K.K., Wan, K.M., \& King, R.B. (2021). Performance over enjoyment? Effect of game-based learning on learning outcome and flow experience. \textit{Frontiers in Education}, 6, 660376.

\bibitem{ref8}
Karran, A.J., Demazure, T., \& Hudelot, C. (2022). Designing for confidence: The impact of visualizing artificial intelligence decisions. \textit{Frontiers in Neuroscience}, 16.

% --- Can-specific (Kategori B) ---

\bibitem{ref9}
Schroeder, N.L., Davis, C.S., \& Yang, Y. (2025). Designing and learning with pedagogical agents: An umbrella review. \textit{Journal of Educational Computing Research}, 62(8).

\bibitem{ref10}
Meng, Z., Fu, Z., Liu, J., Wang, H., \& Pan, Y. (2024). Real-Time Hand Gesture Monitoring Model Based on MediaPipe's Registerable System. \textit{Sensors}, 24(19), 6262. \textit{[Catatan: DEPRECATED post-pivot 16/6/26 --- login gesture diganti nomor absen + superadmin. Dipertahankan untuk dokumentasi historis canvas drawing input.]}

\bibitem{ref11}
Wang, N., \& Lester, J. (2023). K-12 education in the age of AI: A call to action for K-12 AI literacy. \textit{International Journal of Artificial Intelligence in Education}, 33(3), 631--640.

\bibitem{ref12}
Yim, I. H. Y., \& Su, J. (2025). Artificial intelligence literacy education in primary schools: A review. \textit{International Journal of Technology and Design Education}, 35, Article 57.

\bibitem{ref13}
Lee, S., Mott, B., Ottenbreit-Leftwich, A., Scribner, A., Taylor, S., Glazewski, K., \& Lester, J. (2021). AI-infused collaborative inquiry in upper elementary school: A game-based learning approach. \textit{Proceedings of the AAAI Conference on Artificial Intelligence}, 35(17), 15653--15661.

\bibitem{ref14}
Gupta, A., Lee, S., Mott, B., Chakraburty, S., Scribner, A., Glazewski, K., Ottenbreit-Leftwich, A., \& Lester, J. (2024). Supporting upper elementary students in learning AI concepts with story-driven game-based learning. \textit{Proceedings of the AAAI Conference on Artificial Intelligence}, 38(21), 23321--23329.

\bibitem{ref15}
Hsu, T. C., Abelson, H., Lao, N., Tseng, Y. H., \& Lin, Y. T. (2021). Behavioral-pattern exploration and development of an instructional tool for young children to learn AI. \textit{Computers and Education: Artificial Intelligence}, 2, Article 100060.

\bibitem{ref16}
Ocak, C., Kopcha, T. J., \& Dey, R. (2023). An AI-enhanced pattern recognition approach to temporal and spatial analysis of children's embodied interactions. \textit{Computers and Education: Artificial Intelligence}, 4, Article 100155.

\bibitem{ref17}
Shokeen, E., Katirci, N., Williams-Pierce, C., \& Fan, C. (2022). Children learning to sketch: Sketching to learn. \textit{Information and Learning Sciences}, 123(7/8), 482--497.

\bibitem{ref18}
Chen, Z., Duan, Z., Qi, Y., \& Wang, W. (2025). From line sketches to contextualized scenes: An AR scaffold for narrative drawing based on sketch recognition and its motivational effects. \textit{Proceedings of the 2025 ACM/IEEE International Conference on Human-Robot Interaction (HRI)}, 1--10.

\bibitem{ref19}
Wang, C. J., Zhong, H. X., Chiu, P. S., Chang, J. H., \& Chen, M. Y. (2022). Research on the impacts of cognitive style and computational thinking on college students in a visual artificial intelligence course. \textit{Frontiers in Psychology}, 13, Article 864416.

\bibitem{ref20}
Wu, S. Y., \& Su, Y. S. (2021). Visual programming environments and computational thinking performance of fifth-and sixth-grade students. \textit{Journal of Educational Computing Research}, 59(7), 1413--1441.

% --- General/Shared additions (Kategori A baru) ---

\bibitem{ref21}
Yin, M., Wortman Vaughan, J., \& Wallach, H. (2019). Understanding the effect of accuracy on trust in machine learning models. \textit{Proceedings of the 2019 CHI Conference on Human Factors in Computing Systems}, 1--12.

\bibitem{ref22}
Batanero, C., \& \'Alvarez-Arroyo, R. (2023). Teaching and learning of probability. \textit{ZDM -- Mathematics Education}, 56, 5--17.

\bibitem{ref23}
Zhan, Z., Tong, Y., Lan, X., \& Zhong, B. (2024). A systematic literature review of game-based learning in artificial intelligence education. \textit{Interactive Learning Environments}, 32(4), 1189--1212.

\bibitem{ref24}
Liao, Q. V., Gruen, D., \& Miller, S. (2020). Questioning the AI: Informing design practices for explainable AI user experiences. \textit{Proceedings of the 2020 CHI Conference on Human Factors in Computing Systems}, 1--15.

\bibitem{ref25}
Feldman-Maggor, Y., Cukurova, M., Kent, C., Luckin, R., \& Baur-Wolfson, M. (2025). The impact of explainable AI on teachers' trust and acceptance of AI EdTech recommendations. \textit{International Journal of Artificial Intelligence in Education}, 35(1), 1--25.

% --- Methodology (Kategori A existing, methodology) ---

\bibitem{ref26}
Pressman, R.S. (2014). \textit{Software Engineering: A Practitioner's Approach} (8th ed.). New York: McGraw-Hill.

\bibitem{ref27}
Sommerville, I. (2015). \textit{Software Engineering} (10th ed.). Pearson.

\bibitem{ref28}
Sugiyono. (2020). \textit{Metode Penelitian Kuantitatif, Kualitatif, dan R\&D}. Bandung: Alfabeta.

\bibitem{ref29}
Creswell, J.W. (2014). \textit{Research Design: Qualitative, Quantitative, and Mixed Methods Approaches} (4th ed.). Sage Publications.

\bibitem{ref30}
Ishikawa, K. (1986). \textit{Guide to Quality Control} (2nd ed.). Tokyo: Asian Productivity Organization.

\bibitem{ref31}
Brooke, J. (1996). SUS: A `Quick and Dirty' Usability Scale. In \textit{Usability Evaluation in Industry} (pp. 189--194). Taylor \& Francis.

\bibitem{ref32}
Bangor, A., Kortum, P.T., \& Miller, J.T. (2008). An empirical evaluation of the System Usability Scale. \textit{International Journal of Human-Computer Interaction}, 24(6), 574--594.

\end{thebibliography}
